\scanpage{190}
3) Soit $G$ un groupe \deleted{fini} étale sur $S$, et considérons un $G_X$-torseur $P$ sur $X$ (à droite). \deleted{Pour toute} \inserted{Considérons une} représentation de $G$ \deleted{dans} un module qu. coh. $F$ sur $S$. Alors $G_X$ opère sur $F_X$, et on constate qu'il opère en laissant invariant la connexion canonique de $F_X$. Comme les Modules à connexion intégrable \inserted{rel/$S$} sur des objets du site étale de $X$ forment évidemment un champ, on voit que $G_X$ \inserted{(considéré comme \inserted{faisceau en} groupes sur $S$)} opère sur \deleted{\unsure{chaque|cha}} section \inserted{$F_X$,} du dit \deleted{\unsure{$F$|$E$}} champ, qu'on peut donc tordre par le $G_X$-torseur $P$. On trouve donc un module à connexion intégrable
\[
 P\mathbin{\times^{G_X}} F_X\; .
\]
\margin{En fait, on trouve m\textsuperscript{e}, par le m\textsuperscript{e} argument, une \uline{stratification} \uline{intégrable} sur $P\mathbin{\times^{G_X}}F_X$. Cela revient au m\textsuperscript{e} si $X/S$ lisse et $S$ de car. 0.}

Cette façon de construire des \deleted{connexions} \inserted{stratifications} intégrables est \deleted{\unsure{une construction|une connexion} \illegible{1}} universelle dans une large mesure. \margin{\unsure{Jouer|Prouver|Donnons} stabilité par sous-modules, quotients etc} \deleted{Supposons que $X\longrightarrow S$ soit plat, \deleted{\unsure{et loc.}} \inserted{\deleted{\unsure{loc.}} de prés. finie,} à fibres géométriques connexes, et qu'il ait une section $e$. Supposons} Du moins si $S$ est réduit au spectre d'un corps $k$, et que $X$ est p. ex. connexe et \deleted{de} loc. de présentation finie sur $k$, muni d'un pt $x\in X$ \deleted{\unsure{rationnel sur}} \inserted{\unsure{résiduel|rationnel} sur} $k$, on peut former le pro-groupe \inserted{strict} \deleted{\unsure{par}} (pas néc. profini si $X$ n'est pas normal) $\underline{\pi}_1(X,x)$, et la catégorie des Modules \inserted{\uline{coh.}} à stratification intégrables sur $X$ doit être équivalente à celle des représentations linéaires de dim finie de $\underline{\pi}_1(X,x)$. On doit pouvoir généraliser au cas d'une base quelconque $S$, en se
% archive sheet number: 189
